\usepackage{amsfonts}
\usepackage{amssymb}
\usepackage{amsmath}
\usepackage{enumitem} % paquete para listas y cambiar la etiqueta
\usepackage{tabularray} % para tablas
\UseTblrLibrary{amsmath}
\usepackage{bm} % Para poner negrita en fórmulas
\usepackage{siunitx}
\usepackage{microtype}
%%%%%%%%%%%%%%%%%%%%%%%%%%%%%%%%
% Unidades SIUNIX
%%%%%%%%%%%%%%%%%%%%%%%%%%%%%%%%%%%
\sisetup{%
    %alsoload = synchem,
    detect-family,%
    per-mode = fraction,
    %per-mode = symbol-or-fraction,
    output-decimal-marker = {,},
    inter-unit-product = \cdot,
    exponent-product = \cdot,
    retain-explicit-plus % consigo que si pongo + en el número aparezca  
}

\DeclareSIUnit\cv{CV}
\DeclareSIUnit\kilopondio{kp}
\DeclareSIUnit\dina{din}
\DeclareSIUnit\atmospher{atm}
\DeclareSIUnit\calorie{cal}
\DeclareSIUnit\erg{erg}
\DeclareSIUnit\kwh{\kilo\watt\hour}
\DeclareSIUnit\kilojoule{\kilo\joule}
\DeclareSIUnit\uma{u}
%\DeclareSIUnit\mol{mol}
\DeclareSIUnit\mmhg{mmHg}
%\DeclareSIUnit{\molar}{M} está declarado en chemmacros
\DeclareSIUnit{\M}{M}
\usepackage{chemformula}
\usepackage{chemmacros}
\usepackage{elements}
\usepackage{float} % para poner la gráfica donde considere
\usepackage{varwidth} % Se necesita para algunos tcolorbox
\usepackage{hyperref}
\hypersetup{%
    colorlinks=true,
    linkcolor = darkblue
    %urlcolor=myOrange3
}
\usepackage{xparse} 
\usepackage{pdflscape}
\usepackage{setspace} % Para establecer el espacio , se hace con \spacing{1.2}
\usepackage{paracol}
\usepackage{lewis}
\usepackage[normalem]{ulem}
\usepackage{soul}
\usepackage{soulutf8}
\usepackage{chemfig}
\usepackage{booktabs}
\usepackage{comment}
\usepackage{xstring}% Necesario para ifEqcases
\usepackage{cancel}
\usepackage{multicol}
\usepackage[spanish=mexican]{csquotes} %para las comillas
\usepackage{epsfig} %para pasar a eps a pdf

\usepackage{mathpazo}
\usepackage{euler}
%%%%%%%%%%%%%%%%%%%%%%%%%%%%%%%%%%%%%%%%%%%%%%%%%%%%%%%%%%%%%%%%%%%%%%%%%%%%
% TIKZ y librería
%%%%%%%%%%%%%%%%%%%%%%%%%%%%%%%%%%%%%%%%%%%%%%%%%%%%%%%%%%%%%%%%%%%%%%%%%%%%
\usepackage{tikz}
\usetikzlibrary{matrix,arrows, positioning,shadows,shadings,backgrounds,
calc, shapes, tikzmark}

\UseTblrLibrary{siunitx}

%%%%%%%%%%%%%%%%%%%%%%%%%%%%%%%%%%%%%%%%%%%%%%%%%%%%%%%%%%%%%%%%%%%%%%%%%%%%

\usepackage{mathtools} % lo he cargado para que me de acceso a rcases
\usepackage{siunitx} 



\usepackage{subfig} %para poner varias imágenes juntas con subfloat

\usepackage{upgreek}

\usepackage{fontawesome}